
\subsubsection{Uniform Overconfidence Interpretation}

The results reveal a two-layer structure:

\begin{itemize}
\item \textbf{Surface layer:} Category-specific ECE variation exists. Finance/Economics questions are better calibrated (ECE=0.152) than Misconceptions/Myths questions (ECE=0.251).
\item \textbf{Deep layer:} The underlying miscalibration is uniform. All categories require maximum temperature smoothing (T=10.0), suggesting that the confidence-accuracy mapping requires the same correction across all semantic domains.
\end{itemize}

This pattern is consistent with RLHF training inducing global overconfidence that dominates category-specific effects. The model is not differently miscalibrated across categories; it is uniformly overconfident in the same direction with the same severity.

\subsubsection{Implications for Calibration Research}

The negative result on h-m2 indicates that category-specific temperature scaling cannot outperform global temperature scaling when optimal temperatures are uniform. The variation in ECE across categories appears to arise from variation in the difficulty or accuracy of questions within each category rather than from systematic differences in the confidence-accuracy relationship.

\subsubsection{Limitations}

\begin{enumerate}
\item \textbf{Single model:} Only Llama-2-7B was evaluated. Results may differ for other architectures or model scales.
\item \textbf{Temperature bound saturation:} T=10.0 is the upper bound of the search range. True optima may lie beyond T=10, and cluster differentiation could emerge at higher temperatures.
\item \textbf{Simulated h-m1 data:} Due to GPU constraints, confidence distribution analysis (h-m1) used simulated data calibrated to ECE patterns, introducing circularity that limits the independence of that finding.
\item \textbf{NLL vs ECE optimization:} Temperature was optimized using negative log-likelihood rather than ECE directly. Alternative loss functions may yield different results.
\item \textbf{Cluster definition:} The 7 clusters are defined by expert judgment rather than empirical clustering methods.
\end{enumerate}
